
On the same AlphaFold2 inference code and the same input, accelerator
choice alone did not determine what a user actually got. A single TPU
v5e chip beat a T4 GPU by 27.8$\times$ at steady state on our August
baseline, a figure Section~\ref{sec:limitations} qualifies. That
number also described one of eight allocated chips; the other seven sat at
0\,MB until we explicitly asked JAX to use them
(Section~\ref{sec:parallelism}), and the pod's list price reflected
all eight regardless (Section~\ref{sec:cost}). Vectorizing with
\texttt{jax.vmap} never recovered them. Mapping independent queries
across devices with \texttt{jax.pmap} did, giving eight chips
6.53--7.91$\times$ the throughput of one on a matched grid. Automatic sharding, meant to make that mapping
unnecessary, left the per-chip footprint at its single-chip size,
consistent with replication. Our best explanation is that AlphaFold2's
Haiku modules carry no sharding annotations for it to propagate from,
which we could not confirm from retained compiler evidence. Separately,
we showed that an ensemble built outside the model can be mapped
across all eight chips with an explicit \texttt{pmap} and
\texttt{jax.lax.pmean} reduction, without changing AlphaFold2's
source. None of this was
visible from steady-state throughput alone. The cold-start cost
appeared only in a profiler trace (Section~\ref{sec:profiling}), the
baseline drift only when we reran the same benchmark weeks later
(Section~\ref{sec:results-hardware}), and the idle-chip cost only when
the slice was priced as it is billed (Section~\ref{sec:cost}).

\subsection*{Future work}

Three extensions follow directly from what this paper could not test.
Google's TPU 8i, announced before these measurements but not evaluated
here \citep{googletpuv8ann2026}, is inference-specialized silicon.
Testing it against the first-call compilation costs characterized here
is a natural next step once it is available.
FastFold's Dynamic Axial Parallelism is a model-parallel alternative
to the data-parallel approach this paper uses. Whether it transfers to
TPU, and how it compares to the \texttt{pmap}-based approach here, is
untested. And the single sequence length and single TPU generation
noted in Section~\ref{sec:limitations} are the most direct gaps to
close: longer sequences, multimer inputs, and later TPU generations
would each test whether the mechanisms identified here, not just the
specific numbers, generalize beyond this study's setup.
